\documentclass[10pt,a4paper]{amsart}
\usepackage{amsmath,amssymb,amsthm,graphicx}
% Retained preamble for the Stacks Project chapter snapshot in this directory.
% The chapter's own preamble.tex is a copy of the shared Stacks Project preamble; the block
% below is sliced MECHANICALLY, line for line, from it. Nothing was added, removed or
% redefined; the wrapper class, the xr-hyper/hyperref wiring and the externaldocument list
% are omitted because the assembler supplies the document class.
%
% lines taken from preamble.tex: 3, 5, 9, 11, 12, 156, 157, 158, ... (51 lines total)
%
% package coverage restored from the chapter's own preamble (2026-10-13): xy, xy-2cell

\usepackage[all]{xy}
\xyoption{2cell}
\UseAllTwocells

\newenvironment{reference}{\comment}{\endcomment}
\newenvironment{slogan}{\comment}{\endcomment}
\newenvironment{history}{\comment}{\endcomment}
\theoremstyle{plain}
\newtheorem{theorem}[subsection]{Theorem}
\newtheorem{proposition}[subsection]{Proposition}
\newtheorem{lemma}[subsection]{Lemma}
\theoremstyle{definition}
\newtheorem{definition}[subsection]{Definition}
\newtheorem{example}[subsection]{Example}
\newtheorem{exercise}[subsection]{Exercise}
\newtheorem{situation}[subsection]{Situation}
\theoremstyle{remark}
\newtheorem{remark}[subsection]{Remark}
\newtheorem{remarks}[subsection]{Remarks}
\numberwithin{equation}{subsection}
\def\lim{\mathop{\mathrm{lim}}\nolimits}
\def\colim{\mathop{\mathrm{colim}}\nolimits}
\def\Spec{\mathop{\mathrm{Spec}}}
\def\Hom{\mathop{\mathrm{Hom}}\nolimits}
\def\Ext{\mathop{\mathrm{Ext}}\nolimits}
\def\SheafHom{\mathop{\mathcal{H}\!\mathit{om}}\nolimits}
\def\SheafExt{\mathop{\mathcal{E}\!\mathit{xt}}\nolimits}
\def\Sch{\mathit{Sch}}
\def\Mor{\mathop{\mathrm{Mor}}\nolimits}
\def\Ob{\mathop{\mathrm{Ob}}\nolimits}
\def\Sh{\mathop{\mathit{Sh}}\nolimits}
\def\NL{\mathop{N\!L}\nolimits}
\def\CH{\mathop{\mathrm{CH}}\nolimits}
\def\proetale{{pro\text{-}\acute{e}tale}}
\def\etale{{\acute{e}tale}}
\def\QCoh{\mathit{QCoh}}
\def\Ker{\mathop{\mathrm{Ker}}}
\def\Im{\mathop{\mathrm{Im}}}
\def\Coker{\mathop{\mathrm{Coker}}}
\def\Coim{\mathop{\mathrm{Coim}}}
\DeclareMathSymbol{\boxtimes}{\mathbin}{AMSa}{"02}
\def\QCohstack{\mathcal{QC}\!\mathit{oh}}
\def\Cohstack{\mathcal{C}\!\mathit{oh}}
\def\Spacesstack{\mathcal{S}\!\mathit{paces}}
\def\Quotfunctor{\mathrm{Quot}}
\def\Hilbfunctor{\mathrm{Hilb}}
\def\Curvesstack{\mathcal{C}\!\mathit{urves}}
\def\Polarizedstack{\mathcal{P}\!\mathit{olarized}}
\def\Complexesstack{\mathcal{C}\!\mathit{omplexes}}
\def\Pic{\mathop{\mathrm{Pic}}\nolimits}
\def\Picardstack{\mathcal{P}\!\mathit{ic}}
\def\Picardfunctor{\mathrm{Pic}}
\def\Deformationcategory{\mathcal{D}\!\mathit{ef}}
\usepackage{cleveref}
\crefname{originallemma}{original lemma}{original lemmas}
\Crefname{originallemma}{Original lemma}{Original lemmas}
\makeatletter
\newlabel{equation-to-solve}{{08S4}{source}{}{sourcepointer-1}{}}
\newlabel{equation-to-solve@cref}{{[originallemma][0][]08S4}{[1][0][]source}{}{}{}}
\makeatother
\crefname{originallemma}{original lemma}{original lemmas}
\Crefname{originallemma}{Original lemma}{Original lemmas}
\makeatletter
\newlabel{more-algebra-definition-local-complete-intersection}{{07D0}{source}{}{sourcepointer-2}{}}
\newlabel{more-algebra-definition-local-complete-intersection@cref}{{[originallemma][0][]07D0}{[1][0][]source}{}{}{}}
\makeatother
\crefname{originallemma}{original lemma}{original lemmas}
\Crefname{originallemma}{Original lemma}{Original lemmas}
\makeatletter
\newlabel{more-algebra-lemma-conormal-sequence-H1-regular-ideal}{{07CX}{source}{}{sourcepointer-3}{}}
\newlabel{more-algebra-lemma-conormal-sequence-H1-regular-ideal@cref}{{[originallemma][0][]07CX}{[1][0][]source}{}{}{}}
\makeatother
\crefname{originallemma}{original lemma}{original lemmas}
\Crefname{originallemma}{Original lemma}{Original lemmas}
\makeatletter
\newlabel{more-algebra-lemma-quasi-regular-ideal-finite-projective}{{08RK}{source}{}{sourcepointer-4}{}}
\newlabel{more-algebra-lemma-quasi-regular-ideal-finite-projective@cref}{{[originallemma][0][]08RK}{[1][0][]source}{}{}{}}
\makeatother
\crefname{originallemma}{original lemma}{original lemmas}
\Crefname{originallemma}{Original lemma}{Original lemmas}
\makeatletter
\newlabel{more-algebra-section-ideals}{{07CU}{source}{}{sourcepointer-5}{}}
\newlabel{more-algebra-section-ideals@cref}{{[originallemma][0][]07CU}{[1][0][]source}{}{}{}}
\makeatother
\title{Stacks, Deformation Theory: lemma-existence-lci}
\author{The Stacks Project Authors}
\date{stacks-project:defos.tex:local-snapshot-20261012; GNU FDL 1.2 or later}
\begin{document}
\maketitle
\noindent\textit{Editorial note.} Statement span is lines 640-644 of sources/stacks-defos/pinned-original-source.tex; proof span is lines 646-675, the complete author proof (28 body lines).\par\medskip
\noindent\textit{Editorial note.} Navigation stubs: 5 cited-but-not-included labels are supplied as non-load-bearing pointers with their pinned tags where the snapshot tag index has them.\par\medskip
\noindent\textit{Editorial note.} extra\_packages is empty on purpose: the retained preamble already loads xy with options and a second declaration clashes in the assembler wrapper.\par\medskip
\noindent\textit{Editorial note.} Text is retained verbatim, including the author's own typos and placeholders.\par\medskip
\noindent\textit{Editorial note.} Queue index 21 of queue\_defos.json (pb=1261, stmt 640-644, proof 646-675).\par\medskip
\noindent\textit{Allowed standard prerequisites (editorial).}\begin{itemize}
\item Explicit prerequisite boundary. The cited inputs are prerequisites; the argument in the retained span is not replaced by any cited result.
\end{itemize}\medskip
\ifdefined\hypertarget\hypertarget{sourcepointer-1}{}\fi
\noindent\textit{Source reference.} Original lemma 08S4: sources/stacks-defos/pinned-original-source.tex, source lines 260--267. Cited by the retained statement or proof but not part of the collected spans; supplied only as a navigation pointer. The argument does not depend on its content.\par\medskip
\ifdefined\hypertarget\hypertarget{sourcepointer-2}{}\fi
\noindent\textit{Source reference.} Original lemma 07D0: sources/stacks-defos/pinned-original-source.tex, source lines 675--675. Cited by the retained statement or proof but not part of the collected spans; supplied only as a navigation pointer. The argument does not depend on its content.\par\medskip
\ifdefined\hypertarget\hypertarget{sourcepointer-3}{}\fi
\noindent\textit{Source reference.} Original lemma 07CX: sources/stacks-defos/pinned-original-source.tex, source lines 675--675. Cited by the retained statement or proof but not part of the collected spans; supplied only as a navigation pointer. The argument does not depend on its content.\par\medskip
\ifdefined\hypertarget\hypertarget{sourcepointer-4}{}\fi
\noindent\textit{Source reference.} Original lemma 08RK: sources/stacks-defos/pinned-original-source.tex, source lines 675--675. Cited by the retained statement or proof but not part of the collected spans; supplied only as a navigation pointer. The argument does not depend on its content.\par\medskip
\ifdefined\hypertarget\hypertarget{sourcepointer-5}{}\fi
\noindent\textit{Source reference.} Original lemma 07CU: sources/stacks-defos/pinned-original-source.tex, source lines 675--675. Cited by the retained statement or proof but not part of the collected spans; supplied only as a navigation pointer. The argument does not depend on its content.\par\medskip
\begin{lemma}
\label{lemma-existence-lci}
If $A \to B$ is a local complete intersection ring map, then
there exists a solution to (\ref{equation-to-solve}).
\end{lemma}

\begin{proof}[First proof]
Write $B = A[x_1, \ldots, x_n]/J$. By More on Algebra, Definition
\ref{more-algebra-definition-local-complete-intersection}
the ideal $J$ is Koszul-regular. This implies $J$ is $H_1$-regular and
quasi-regular, see More on Algebra, Section \ref{more-algebra-section-ideals}.
Let $J' \subset A'[x_1, \ldots, x_n]$
be the inverse image of $J$. Denote $I[x_1, \ldots, x_n]$ the
kernel of $A'[x_1, \ldots, x_n] \to A[x_1, \ldots, x_n]$.
By More on Algebra, Lemma
\ref{more-algebra-lemma-conormal-sequence-H1-regular-ideal} we have
$I[x_1, \ldots, x_n] \cap (J')^2 = J'I[x_1, \ldots, x_n] =
JI[x_1, \ldots, x_n]$. Hence we obtain a short exact sequence
$$
0 \to I \otimes_A B \to J'/(J')^2 \to J/J^2 \to 0
$$
Since $J/J^2$ is projective (More on Algebra, Lemma
\ref{more-algebra-lemma-quasi-regular-ideal-finite-projective})
we can choose a splitting of this sequence
$$
J'/(J')^2 = I \otimes_A B \oplus J/J^2 
$$
Let $(J')^2 \subset J'' \subset J'$ be the elements which map to the
second summand in the decomposition above. Then
$$
0 \to I \otimes_A B \to A'[x_1, \ldots, x_n]/J'' \to B \to 0
$$
is a solution to (\ref{equation-to-solve}) with $N = I \otimes_A B$.
The general case is obtained by doing a pushout along the given
map $I \otimes_A B \to N$.
\end{proof}

\end{document}
